\documentclass[12pt]{article}
\usepackage[utf8]{inputenc}
\usepackage[T1]{fontenc}
\usepackage{hyperref}
\title{Pashto Language Resources for K-12 Curriculum Development}
\author{Compiled by Claude Sonnet 5.5 (claude-sonnet-5-5) for Andr\'as Kornai}
\date{}
\begin{document}
\maketitle

\section{Introduction}
Pashto (ISO 639-3 \texttt{pus}) is an Eastern Iranian language of the Indo-European family, spoken in Afghanistan and Pakistan, with estimates of 40 to 60 million speakers. It is an official language of Afghanistan and the second-largest provincial language of Pakistan, and is written in a Perso-Arabic script with additional letters. This dossier lists only what turned up in a small number of web searches; it is not exhaustive, and the searches found little on state school curricula.

\section{Primary Educational Resources}
\begin{itemize}
\item German publisher wardak.de describes a set of Pashto teaching materials for primary-level learners (Pashto as a second language, reading books, exercise book, grammar book, teacher's guide); the listings were inconsistent about page counts and completion, so verify with the publisher: \url{https://wardak.de/articles/pashto-afghanisch-als-zweite-sprache/}
\item \emph{Pashto: An Elementary Textbook} (Inomkhojayev), a college-level two-volume series with audio/video, useful as a reference for adult or secondary second-language teaching rather than primary school: \url{https://utpdistribution.com/9781589017733/pashto}
\end{itemize}
No official Afghan or Pakistani Pashto school textbooks were found in these searches.

\section{Linguistic and Reference Materials}
\begin{itemize}
\item Wikipedia overview of Pashto (status, speakers, script, dialects, references, including Penzl and Sloan, \emph{A Grammar of Pashto}): \url{https://en.wikipedia.org/wiki/Pashto}
\item EBSCO research starter on the Pashto language: \url{https://ebsco.com/research-starters/language-and-linguistics/pashto-language}
\end{itemize}
No dictionary was located in these searches.

\section{Religious and Community Resources}
\begin{itemize}
\item 4training.net Pashto page, noting that the New Testament and major Old Testament books exist in Pashto: \url{https://4training.net/Pashto}
\end{itemize}

\section{Digital and Audio Resources}
\begin{itemize}
\item Global Recordings Network, Southern Pashto: \url{https://www.globalrecordings.net/en/language/pbt}
\item Global Recordings Network, Northern Pashto: \url{https://www.globalrecordings.net/en/language/pbu}
\item Global Recordings Network, Pushto (pus): \url{https://globalrecordings.net/en/language/pus}
\item ``Learn Pashto'' mobile app (phrases with audio): \url{https://apps.apple.com/us/app/id1176460224}
\end{itemize}

\section{Curriculum Development Context}
Pashto has a large speaker base and official status, but the searches did not surface a native-language K-12 curriculum or textbook series. Variation between Northern and Southern varieties should be considered when choosing a standard. Further work should consult Afghan and Pakistani ministry textbooks directly, which were not found here.

\end{document}